\section{Cloud Native}
Der Begriff \textbf{Cloud Native} wurde 2013 von Matt Stine geprägt und unter anderem durch Netflix auf einem AWS re:Invent Talk bekannt gemacht. Es gibt keine einheitliche Definition für Cloud-Native-Architektur. Grundsätzlich soll Cloud Native dabei helfen, Webanwendungen zu entwickeln, die \emph{skalierbar} und \emph{hochverfügbar} sind. Gleichzeitig soll durch Cloud Native \emph{agil} entwickelt werden, also eine schnellstmögliche Feature-Anpassung mit möglichst kleinem Risiko. Dabei helfen \textit{Microservices} statt einer monolithischen Architektur, \textit{Container}, \textit{DevOps} sowie \textit{CI/CD-Pipelines}. Viele Lösungen wie Komponenten, \textit{Logging} und \textit{Tracing} sind dabei standardisiert.
\begin{figure}[H]
    \centering
    \includegraphics[width=\textwidth]{CloudNativeSaulen.png}
    \caption{Die 4 Saulen von Cloud Native}
\end{figure}

\begin{figure}[H]
    \centering
    \includegraphics[width=\textwidth]{CloudNativeMap.png}
    \caption{Cloud Native Landscape (CNCF)}
\end{figure}

\subsection{DevOps}
Kombination aus \textit{Development} und \textit{Operations}.
Entwicklung und Betrieb arbeiten \emph{zusammen statt getrennt}.
\textbf{Kern-Ideen:} Automatisierung von Builds, Tests und Deployments,
schnelle Feedbackschleifen (kurze Release-Zyklen) und gemeinsame
Verantwortung für Code und Infrastruktur.

\subsection{CI/CD -- Continuous Integration / Continuous Delivery}
\textbf{CI:} Jeder Code-Push wird automatisch gebaut und getestet,
Fehler werden sofort sichtbar.
\textbf{CD:} Nach erfolgreichem CI wird der Build automatisch
in eine Umgebung ausgeliefert (Staging oder Produktion).

\subsection{Microservices}
Architekturstil, bei dem eine Anwendung in viele kleine, unabhängige
Dienste aufgeteilt wird. Jeder Service hat eine \emph{klar abgegrenzte Aufgabe},
kommuniziert über \textit{APIs} (REST oder Message Queues) und kann separat
deployt und skaliert werden. Gegenteil: \textit{Monolith}.

\subsection{Monolith}
Eine monolithische Architektur bedeutet, dass alle Elemente einer Software in einer
zusammenhängenden Codebasis liegen. Die grafische Benutzeroberfläche, die Geschäftslogik
und der Datenzugriff sind eng miteinander verbunden. Dies führt dazu, dass die Codebasis
stark anwächst, wodurch Fehler schwerer nachvollziehbar werden. Bereits eine kleine
Änderung erfordert, dass die gesamte Anwendung neu deployt wird. Die beteiligten Teams
sind häufig groß und müssen bereichsübergreifend zusammenarbeiten.

\begin{figure}[H]
    \centering
    \includegraphics[width=\textwidth]{CloudNativeBesipel.png}
    \caption{Microservices vs.\ Monolith}
\end{figure}
